\begin{titlepage}\thispagestyle{empty}
\begin{tikzpicture}[remember picture,overlay]
  \begin{scope}
    \clip(current page.north west) rectangle ([yshift=-76mm]current page.north east);
    \node[inner sep=0pt,anchor=center] at ([yshift=-38mm]current page.north)
      {\includegraphics[width=\paperwidth]{assets/cover_network.png}};
  \end{scope}
  \fill[CoverTeal]([yshift=-76mm]current page.north west)
    rectangle ([yshift=-88mm]current page.north east);
  \node[anchor=north west,inner sep=0pt,text=black] at ([xshift=18mm,yshift=-109mm]current page.north west)
    {\fontsize{32}{40}\selectfont\sffamily\bfseries Python 网络爬虫};
  \node[anchor=north west,inner sep=0pt,text=black] at ([xshift=19mm,yshift=-129mm]current page.north west)
    {\fontsize{23}{30}\selectfont\sffamily\bfseries 从请求到数据};
  \node[anchor=north west,inner sep=0pt,text=Ink] at ([xshift=20mm,yshift=-151mm]current page.north west)
    {\fontsize{18}{23}\selectfont\bfseries Python Web Scraping Notes};
  \node[anchor=north west,inner sep=0pt,text=black!58,align=left] at ([xshift=23mm,yshift=-177mm]current page.north west)
    {\fontsize{12}{24}\selectfont
      \begin{tabular}{@{}l@{\;}l@{}}
      作者：&雪落长安无声\\
      日期：&2026.11.05\\
      版本：&2026 年修订
      \end{tabular}};
  \begin{scope}[shift={([xshift=-39mm,yshift=-207mm]current page.north east)}]
    \foreach \x/\y in {-1.2/0.4,-0.5/1.1,0.5/1.4,1.3/0.5,0.7/-0.5,-0.6/-0.8}{
      \draw[CoverGold!65,line width=.5pt] (0,0)--(\x,\y);
      \fill[CoverGold] (\x,\y) circle (2.2pt);
    }
    \fill[Teal] (0,0) circle (5pt);
  \end{scope}
  \node[anchor=center,text=black,align=center] at ([yshift=37mm]current page.south)
    {\fontsize{12}{19}\selectfont\itshape 从网页到数据，每一步都可解释、可复现。};
\end{tikzpicture}
\null
\end{titlepage}
\thispagestyle{empty}\vspace*{32mm}
\begin{center}
{\LARGE\bfseries Python 网络爬虫：从请求到数据}\par\vspace{15mm}
原理、解析与工程实践\par\vspace{22mm}
{\large 雪落长安无声}\par\vspace{18mm}
请求 · 解析 · 清洗 · 存储\par\vspace{12mm}
2026.11.05
\end{center}\vfill
从可复现的小样本出发，依次练习请求、解析、存储与排错。代码示例与正文分开保存，便于逐行运行和核对。
